% !PRINTED 574
\weberpage{574}
%<<CONT
in which the congruence is referred to the periods, and thus (\S\ 152, 1.)
\weq{2}{\begin{array}{rll}
\text{in case a),} & \text{when } \nu \equiv \pm\nu', & \\
\text{b),} & \text{when } \nu \equiv \pm\nu', \ \pm i\nu' & (\mathrm{mod}\ \mu), \\
\text{c),} & \text{when } \nu \equiv \pm\nu', \ \pm\varrho\nu', \ \pm\varrho^2\nu'. &
\end{array}}
Now let
\weq{3}{\mu = \mathfrak{a}\,\mathfrak{m}, \quad \mu_1 = \mathfrak{a}_1\,\mathfrak{m}}
be two integers with the greatest common ideal divisor $\mathfrak{m}$ in the field $\Omega$. Then $T_\mu$ and $T_{\mu_1}$ have a common linear divisor if
\weq{4}{\tau\left(\frac{\nu}{\mu}\right) = \tau\left(\frac{\nu_1}{\mu_1}\right),}
and by (2) this takes place then and only then, when
\weq{5}{\mu\nu_1 \equiv \varepsilon\mu_1\nu \quad (\mathrm{mod}\ \mu\mu_1),}
in which $\varepsilon$ is a \gesperrt{unit} of the field $\Omega$, and thus $\varepsilon = \pm 1$ in general, $\varepsilon = \pm 1, \pm i$ when $\Delta = -4$, and $\varepsilon = \pm 1, \pm\varrho, \pm\varrho^2$ when $\Delta = -3$. But from (5) it follows that $\nu$ must be divisible by $\mathfrak{a}$, and $\nu_1$ by $\mathfrak{a}_1$.

If we think of $\nu$ as given and $\nu_1$ as sought, then the congruence (5) is possible only if $\mu_1\nu$ is divisible by $\mu$; and since $\mathfrak{a}$ and $\mathfrak{a}_1$ are relatively prime, $\nu$ must be divisible by $\mathfrak{a}$. If $\alpha$ is an integer in $\Omega$ divisible by $\mathfrak{a}$, and so constituted that $\alpha : \mathfrak{a}$ is relatively prime to $\mathfrak{m}$, then we may accordingly put
\weq{6}{\nu \equiv \alpha\xi \quad (\mathrm{mod}\ \mu)}
(Bd. II, \S\ 166, 7.), and we obtain all the values $\tau\left(\dfrac{\nu}{\mu}\right)$ that are different from one another, and each of them two, four, or six times, if we let $\xi$ run through a complete system of residues with respect to the modulus $\mathfrak{m}$ (with exclusion of zero). One then obtains all the common roots of $T_\mu$ and $T_{\mu_1}$ in the form
\weq{7}{\tau_\xi = \tau\left(\frac{\alpha\xi}{\mu}\right).}
The greatest common divisor $D_{\mathfrak m}$ of $T_\mu$ and $T_{\mu_1}$ therefore has all the quantities (7) as roots. From $D_{\mathfrak m}$ a divisor $T_{\mathfrak m}$ can still be separated off by rational means, one that has as roots only those among the quantities (7) in which $\xi$ is relatively prime to $\mathfrak{m}$, if one frees $D_{\mathfrak m}$ of all the factors that it has in common with a $D_{\mathfrak m'}$, where $\mathfrak{m}'$ is a divisor of $\mathfrak{m}$. Since in this
%CONT>>
